
\subsubsection{2.1 Benchmark Evaluation}

TruthfulQA \citep{Lin2021TruthfulQA} measures whether models produce truthful answers to questions that elicit common human misconceptions. The benchmark reports aggregate accuracy but does not categorize individual failures by error type. FEVER \citep{Thorne2018FEVER} evaluates claim verification against a Wikipedia-derived knowledge base. These benchmarks serve measurement functions but do not provide diagnostic classification of failure modes.

\subsubsection{2.2 Interpretability Methods}

Attention visualization has been used to analyze transformer models. Vig and Belinkov (2019) visualize attention flow in BERT to trace information paths. Clark et al. (2019) identify attention heads that specialize in linguistic phenomena such as coreference resolution. These methods require manual example selection and typically analyze 10-20 cases per study. This small-scale approach limits systematic discovery of failure patterns across benchmark datasets.

The present study automates attention extraction and entropy calculation across 73 benchmark failures, eliminating manual example selection and enabling statistical validation of attention pattern signatures.

\subsubsection{2.3 Correction Methods}

Retrieval-Augmented Generation (RAG; Lewis et al., 2020) addresses factual knowledge gaps by retrieving documents from external corpora and incorporating them into generation context. Chain-of-Thought prompting (COT; Wei et al., 2022) elicits intermediate reasoning steps to improve performance on multi-step tasks.

These methods are typically applied uniformly to all failures without failure-type-specific routing. The present framework introduces failure-type-aware routing based on automated diagnostic classification. However, correction effectiveness is not validated in real-world settings; only synthetic structure validation was performed.
